% NOT PART OF THE PAPER.
%
% Two impossibility results, cut because neither is needed to prove that the
% link exists: they answer questions about what CANNOT work, which is adjacent
% to the problem rather than part of it. Both are complete and proved. They are
% kept here, and listed in plan.md section 8, for a possible separate note.
%
%   Theorem (convex planar). No tetrahedrally symmetric four-component link
%   all of whose components are convex planar satisfies condition (i).
%
%   Theorem (minimal crossings). No tetrahedrally symmetric diagram exists in
%   which the three components at each vertex cross pairwise exactly once and
%   nowhere else -- the twelve-crossing model clasps all six pairs.
%
% Supporting computation: vertex-weave.mjs.

\section{Two obstructions}\label{sec:obstructions}

Both obstructions in this section say that an obvious economy is unavailable.
They are what force the construction of Section~\ref{sec:construction} to look
as elaborate as it does.

\subsection{The components cannot be flat and convex}

Call a closed curve \emph{convex planar} if it lies in a plane and bounds a
convex region there; equivalently, no straight line meets it in more than two
points. Triangles and ellipses are convex planar; a star-shaped polygon is not.

\begin{theorem}\label{thm:convex}
No $\Tet$-symmetric four-component link all of whose components are convex
planar satisfies condition~(i).
\end{theorem}

\begin{proof}
Let $a_i$ be the axis through $v_i$, so that $K_i$ is invariant under the
rotation $r_i$ of order $3$ about $a_i$.

\emph{Step 1: the planes are forced.} A planar simple closed curve spans a
unique plane $P_i$, and $r_i(K_i)=K_i$ forces $r_i(P_i)=P_i$. A rotation through
$120^\circ$ about $a_i$ carries a plane containing $a_i$ to a different such
plane, since $120^\circ\not\equiv0\pmod{180^\circ}$; hence $P_i\perp a_i$.

\emph{Step 2: the half-turn reverses the line of intersection.} For $i\neq j$
the axes $a_i,a_j$ meet at an angle $\arccos(-1/3)$, so $P_i$ and $P_j$ are not
parallel and $\ell=P_i\cap P_j$ is a line. There is a unique element $u\in\Tet$
of order $2$ with $u(K_i)=K_j$ and $u(K_j)=K_i$, namely the half-turn realising
the double transposition $(ij)(kl)$; let $b$ be its axis. Then $u(\ell)=\ell$. A
half-turn carries a line to itself only if the line is $b$, or meets $b$
orthogonally. The first is impossible, as $b$ lies in neither $P_i$ nor $P_j$;
in the second case $u|_\ell$ is the point reflection of $\ell$ in the point
$\ell\cap b$. Concretely, for $i=1$, $j=2$ we have $b$ the $x$-axis, and with
$P_1=\{x+y+z=c\}$ and $P_2=\{x-y-z=c\}$ one computes
$\ell=\{(c,s,-s):s\in\RR\}$, whose direction $(0,1,-1)$ is orthogonal to $b$ and
which meets $b$ at $(c,0,0)$.

\emph{Step 3: chords cannot nest.} Being convex, $K_i$ meets $\ell$ in at most
two points, so (generically) $K_i\cap\ell$ is the pair of endpoints of an
interval $I_i\subset\ell$, and likewise for $K_j$. Two disjoint convex planar
curves lying in distinct planes are linked precisely when $I_i$ and $I_j$
interleave, and unlinked precisely when the intervals nest or are disjoint. By
Step 2, $I_j$ is the image of $I_i$ under the point reflection in
$m=\ell\cap b$. Taking $m$ as the origin of a coordinate on $\ell$ and writing
$I_i=[\alpha,\beta]$, we get $I_j=[-\beta,-\alpha]$. For $I_i\subseteq I_j$ we
would need $-\beta\le\alpha$ and $\beta\le-\alpha$, that is
$\alpha+\beta\ge0$ and $\alpha+\beta\le0$; so $\alpha+\beta=0$, $I_i=I_j$, and
the two curves meet. Nesting is therefore unavailable
(Figure~\ref{fig:intervals}).

\emph{Step 4: neither survivor works.} The group $A_4$ is $2$-transitive on four
points, hence transitive on the six unordered pairs, so every pair falls into
the same one of the remaining two cases. If all pairs interleave then
$\lk(K_i,K_j)=\pm1$ for every pair, contradicting Lemma~\ref{lem:pairs}. If all
pairs are disjoint, then the convex disc $D_i\subset P_i$ bounded by $K_i$ meets
$P_j$ only in the filled interval $I_i$, which is disjoint from $K_j$; so $D_i$
is a disc bounded by $K_i$ and disjoint from every other component, every
three-component sublink is split, and no such sublink is $\Bor$.
\end{proof}

\begin{figure}[htbp]
\centering
\begin{tikzpicture}[scale=0.92,font=\small]
  \foreach \y/\lab in {0/{disjoint}, -1.5/{interleaved}, -3.0/{nested}}
    {\draw[rulegrey] (-3.6,\y) -- (3.6,\y);
     \draw[rulegrey,line width=1pt] (0,\y-0.17) -- (0,\y+0.17);
     \node[rulegrey,font=\scriptsize] at (0,\y+0.38) {$m$};
     \node[anchor=west,font=\scriptsize] at (4.0,\y) {\lab};}
  % disjoint
  \draw[ringB,line width=3pt,line cap=round] (-3.0,0)--(-1.3,0);
  \draw[ringC,line width=3pt,line cap=round] (1.3,0)--(3.0,0);
  % interleaved
  \draw[ringB,line width=3pt,line cap=round] (-3.0,-1.5)--(1.3,-1.5);
  \draw[ringC,line width=3pt,line cap=round] (-1.3,-1.5)--(3.0,-1.5);
  % nested -- greyed, because it is the case the symmetry forbids
  \draw[ringB!35,line width=3pt,line cap=round] (-2.4,-3.0)--(2.4,-3.0);
  \draw[ringC!35,line width=3pt,line cap=round] (-1.1,-3.0)--(1.1,-3.0);
  \node[font=\scriptsize,anchor=west] at (4.0,-3.34) {\textcolor{rulegrey}{(impossible here)}};
\end{tikzpicture}
\caption{The three ways two intervals can sit on the line $\ell=P_i\cap P_j$.
The half-turn exchanging the two components reflects $\ell$ in the point $m$,
so one interval is the mirror image of the other through $m$. Disjoint and
interleaved configurations are compatible with that; nesting is not, and nesting
is the only one of the three that yields a Borromean triple.}
\label{fig:intervals}
\end{figure}

\subsection{The components cannot cross the minimum number of times}

Theorem~\ref{thm:convex} forces the components off their planes. The next
result forces them to cross each other more often than one might hope.

Suppose, as the picture in Figure~\ref{fig:tetra} suggests, that each component
runs along the boundary of its face. Two faces $F_i,F_j$ share an edge whose
endpoints are the two vertices $v_k,v_l$ with $\{k,l\}=\{1,2,3,4\}\setminus\{i,j\}$,
so the components $K_i$ and $K_j$ come close together near $v_k$ and near $v_l$
and nowhere else. The most economical diagram consistent with that picture is
one in which the three components at each vertex cross each other once apiece,
giving $4\times3=12$ crossings in all and exactly two crossings per pair --- the
fewest any link satisfying our conditions could have, since a pair crossing
fewer than twice has nonzero linking number or does not interact at all.

\begin{theorem}\label{thm:minimal}
There is no $\Tet$-symmetric diagram of a four-component link in which, for each
vertex $v_k$, the three components $\{K_i:i\neq k\}$ pairwise cross exactly once
near $v_k$, and no crossings occur elsewhere.
\end{theorem}

\begin{proof}
Fix such a diagram. Near $v_k$ the three components cross pairwise once, so the
relation ``$K_i$ passes over $K_j$ near $v_k$'' is a tournament on three
elements. The stabiliser of $v_k$ in $\Tet$ is $\langle r_k\rangle$, which
permutes those three components in a $3$-cycle and preserves the relation; a
$3$-cycle preserves no linear order on three elements, so the relation is a
cyclic order, and there are exactly two of them. Since $\Tet$ is transitive on
the vertices of $\Delta$, the choice at one vertex determines the choice at all
four, so there are exactly two diagrams to consider, interchanged by reversing
every cyclic order.

Each pair $\{K_i,K_j\}$ crosses exactly twice, at $v_k$ and at $v_l$. Two
crossings between distinct components contribute $\pm1$ each to
$2\lk(K_i,K_j)$, and the two contributions cancel exactly when the \emph{same}
component passes over at both crossings. By Lemma~\ref{lem:pairs} we therefore
need the over-strand of the pair $\{K_i,K_j\}$ to be the same at $v_k$ and at
$v_l$.

Take the cyclic order at $v_k$ to be the orbit of $r_k$, the rotation through
$+120^\circ$ about the axis through $v_k$. In the coordinates
of~\eqref{eq:r1}--\eqref{eq:halfturns} these rotations act on indices as
\[
  r_1=(2\,3\,4),\qquad r_2=(1\,4\,3),\qquad r_3=(1\,2\,4),\qquad r_4=(1\,3\,2),
\]
and the over-strand of a pair at a vertex is read off from the corresponding
cyclic order. The six pairs give
\begin{center}\small
\begin{tabular}{@{}lccc@{}}
\toprule
pair & vertices & over at the first & over at the second\\
\midrule
$\{K_1,K_2\}$ & $v_3,v_4$ & $K_1$ & $K_2$\\
$\{K_1,K_3\}$ & $v_2,v_4$ & $K_3$ & $K_1$\\
$\{K_1,K_4\}$ & $v_2,v_3$ & $K_1$ & $K_4$\\
$\{K_2,K_3\}$ & $v_1,v_4$ & $K_2$ & $K_3$\\
$\{K_2,K_4\}$ & $v_1,v_3$ & $K_4$ & $K_2$\\
$\{K_3,K_4\}$ & $v_1,v_2$ & $K_3$ & $K_4$\\
\bottomrule
\end{tabular}
\end{center}
In every row the over-strand changes, so every pair has linking number $\pm1$.
Reversing all four cyclic orders exchanges the two columns and changes nothing.
Both diagrams are therefore excluded by Lemma~\ref{lem:pairs}.
\end{proof}

Theorems~\ref{thm:convex} and~\ref{thm:minimal} together explain the shape of
what follows: the components must leave their planes, and they must meet each
other more than twice per pair.

